\section{Conclusion and Limitations}
\label{sec:conclusion}

For grades unseen during adaptation, fixing the score and fixing the order are different problems.
Class-bias calibration changes how often a missing grade is predicted but provably cannot change its image ranking against a fixed seen grade.
Blending toward the zero-shot model moves that order only along a single path.
What the missing grade lacks is an adaptation of its own.
TIE supplies one from the grade order: each missing grade inherits the image-dependent displacement of its seen neighbors while keeping its own zero-shot reference, without training and without touching seen-grade scores.
Across four benchmarks and nine adapter configurations, TIE lowers calibrated missing-grade and overall error, improves missing--seen ranking in 34 of 36 adapter--dataset pairs, and on three benchmarks surpasses the best evaluated global Blend.
More broadly, when a label space is ordered, a class missing from adaptation need not be treated as a calibration problem: its neighbors' adaptation can be transferred.

\noindent\textbf{Limitations.}
TIE assumes that adjacent grades share visual evidence; where appearance changes abruptly between grades, or the zero-shot reference is weak, inheritance can help less.
At an endpoint grade, TIE recovers, but cannot exceed, the zero-shot order with the nearest seen grade.
Better ranking does not always yield lower calibrated error, and our missing grades come from controlled masks rather than naturally rare grades.
Learning when and from where to inherit displacement is a natural next step.
